% =====================================================================
%  Stat 236 Final Project Report -- Journal-style template
%  Course guidance: 12pt, single spacing, main body <= ~8 pages
%  (appendix not counted). No code in the main text.
% =====================================================================
\documentclass[12pt]{article}

% ---------- Page layout & fonts ----------
\usepackage[letterpaper,margin=1in,headheight=15pt]{geometry}
\usepackage[T1]{fontenc}
\usepackage[utf8]{inputenc}
\usepackage{lmodern}
\usepackage{microtype}

% ---------- Math, figures, tables ----------
\usepackage{amsmath,amssymb,amsthm,bm}
\usepackage{graphicx}
\usepackage{booktabs}
\usepackage{multirow}
\usepackage{subcaption}
\usepackage{float}
\usepackage{tabularx}
\usepackage[font=small,labelfont=bf,skip=4pt]{caption}

% ---------- Compactness ----------
\usepackage[compact]{titlesec}
\titleformat{\section}{\large\bfseries}{\thesection}{0.6em}{}
\titleformat{\subsection}{\normalsize\bfseries}{\thesubsection}{0.6em}{}
\titlespacing*{\section}{0pt}{1.0ex plus .3ex}{0.5ex}
\titlespacing*{\subsection}{0pt}{0.8ex plus .2ex}{0.3ex}
\usepackage{enumitem}
\setlist{nosep,leftmargin=*}
\setlength{\parskip}{0.3em}
\setlength{\abovecaptionskip}{4pt}
\setlength{\textfloatsep}{8pt plus 2pt minus 2pt}

% ---------- Citations & links ----------
\usepackage[authoryear,round]{natbib}
\bibliographystyle{abbrvnat}
\usepackage[colorlinks=true,linkcolor=blue!50!black,citecolor=blue!50!black,urlcolor=blue!50!black]{hyperref}
\usepackage[nameinlink,capitalise]{cleveref}

% ---------- Placeholder text & notes ----------
\usepackage{lipsum}          % \lipsum[1] etc. -- delete when writing
\usepackage{xcolor}
% Section-owner tags: set \showownerstrue -> \showownersfalse to hide all before submitting
\newif\ifshowowners \showownerstrue
\newcommand{\owner}[1]{\ifshowowners\marginpar{\raggedright\footnotesize\color{red!70!black}\textsf{#1}}\fi}
\newcommand{\todo}[1]{\ifshowowners{\color{red!70!black}\textsf{[TODO: #1]}}\fi}

% ---------- Running header ----------
\usepackage{fancyhdr}
\pagestyle{fancy}
\fancyhf{}
\fancyhead[L]{\small\textit{Stat 236 Final Project}}
\fancyhead[R]{\small\textit{Nicov, Bamzai-Wokhlu, Wilk \& Yang}}
\fancyfoot[C]{\small\thepage}
\renewcommand{\headrulewidth}{0.4pt}

% ---------- Handy macros ----------
\newcommand{\R}{\mathbb{R}}
\newcommand{\E}{\mathbb{E}}
\DeclareMathOperator*{\argmax}{arg\,max}
\newtheorem{definition}{Definition}

% =====================================================================
\begin{document}

% ---------- Journal-style title block ----------
\thispagestyle{plain}
\begin{center}
  {\small\textsc{Stat 236: Network and Text Analysis \quad$\cdot$\quad Fall 2026}}\\[0.6em]
  \rule{\textwidth}{0.8pt}\\[0.5em]
  {\bfseries Contested Classrooms: Text and Network Analysis of Local School Board Meetings}\\[0.6em]
  \rule{\textwidth}{0.4pt}\\[0.8em]
  {\large Stefan Nicov \quad Himal Bamzai-Wokhlu \quad Calista Wilk \quad Lillian Yang}\\[0.3em]
  {\small Harvard University}\\[0.3em]
  {\small Project webpage: \url{https://example.com} \quad$\cdot$\quad Code: \url{https://github.com/himalbw/SchoolPolitics}}
\end{center}

\begin{center}
\begin{minipage}{0.9\textwidth}
\small
\noindent\textbf{Abstract.} We perform text analysis of a
newly released corpus of school board meetings from 2018--2023 to characterize
trends in education politics surrounding the COVID period. Using 80,063 meeting
transcripts from 1,485 districts, we combine bag-of-words measures of partisan
language, topic models, Higher Criticism tests of vocabulary change and
LLM-scored stance. We find that conflict and contested topics such as critical
race theory, book content and parental rights rose sharply during the pandemic,
but they rose in Republican- and Democratic-leaning districts alike, and in counties with varying experiences with school closures.  Partisan
differences in language grew only slightly, peaked in 2021, and mostly reverted
by 2023. The largest vocabulary divide is between professional and local talk
rather than over culture-war issues. The clearest partisan split appears in the
positions argued during public comment in late 2021, and it had closed by 2024.

\vspace{20em}
\end{minipage}
\end{center}
\vspace{0.3em}

\newpage

\tableofcontents

\newpage 

% =====================================================================
\section{Introduction}\label{sec:intro}
Local education politics drew
national attention during the COVID period and recovery. Disputes surrounding school closures, mask mandates, and critical race theory (CRT) transformed school board meetings into heated flashpoints representative of the broader culture war
\citep{onpoint_school_boards_2021,cnn_school_boards_2021,pollock_conflict_2022}. In this paper, we leverage a database of school board meeting transcripts along
with text-analysis methods to quantitatively characterize the politicization of
education surrounding COVID. 

Our dataset DistrictView
\citep{districtview_2026,holman_et_al_2026} is a collection of YouTube-captioned
school board meeting transcripts, which after screening spans 80,063 meetings from
1,485 districts between 2018 and 2023. Figure~\ref{fig:workflow} summarizes our
workflow. To pre-process the transcripts, we restrict to 2018--2023 and to regular
school board meetings, tokenize, and drop stop words, and we append county-level
partisanship from presidential vote shares as covariates. We then pursue two
complementary representations of the text. In the first, we build a bag-of-words
representation, reweighted by TF-IDF, and count mentions of a hand-built dictionary
of 17 themes: four routine school-business controls and 13 political and contested
themes, including two measures of conflict. These themes also guide the LLM
analysis: they locate the contested passages we score, and the attributes that large
language models score through GABRIEL \citep{gabriel} mirror them. In the second, we fit LDA topic models and apply higher criticism
to test whether the vocabulary of meetings changed between the pre-pandemic and
post-pandemic periods.

We present four main findings. First, meeting language
became more partisan post-pandemic (2021 but cooled by 2023. Second, conflict-driven discussions -- such as mentions of critical race theory, book content and
parental rights rose in both groups without differing detectably between them. (CHECK THIS)? 
Third, the trends of topics discussed underwent the most change during COVID, corresponding to a mild temporary polarization between political parties during this period.
Fourth, the sharpest divergence came in the positions argued rather than the words used, as LLM-scored conservative positions rose more in Republican-leaning counties in the second half of 2021, mostly in public comment, before the gap closed by 2024.

\section{Data and Pre-Processing}\label{sec:data}

We use DistrictView, a collection of U.S. school-governance
meeting transcripts from YouTube
\citep{districtview_2026,holman_et_al_2026}.
Our word-based analyses cover 2018--2023.

For the main bag-of-words analysis, we remove unavailable
transcripts, transcripts shorter than 1,000 words under its
tokenization rules, exact duplicates, and recordings that fail
the school-governance eligibility screen. The screen excludes
meetings of other public bodies, joint-body meetings, possible
clips or campaign events, and unresolved cases. Retained
recordings can include regular meetings, special sessions,
workshops, and committees. Text is lowercased and converted
to word tokens after removing timestamps, bracketed annotations,
and punctuation; numerical tokens are excluded from the
bag-of-words features.

The resulting corpus contains 80,063 recordings from 1,485
districts in 49 states, totaling approximately 1.13 billion words
and 140,000 hours of video. The median recording contains
approximately 12,100 words and lasts 1.5 hours.
\Cref{tab:corpus} summarizes this corpus. HC uses the same raw
data and eligibility rules but a separate preprocessing procedure
and a 200-word minimum, producing 83,369 eligible recordings
before the matched-panel restrictions in \cref{sec:m-lda}.
Appendix~\ref{app:cleaning} documents both sample definitions.

\section{Methods}

\subsection{Notation}

We extend standard notation as discussed in class. Let $\mathcal{D} = \{d_1, \dots, d_N\}$ denote the corpus, where each document $d_n$ is the
cleaned auto-caption transcript of one school-board meeting video, and $N = 80{,}063$ is the number of documents.  Each document is a sequence of tokens
\[
  d_n = (w_{n1}, w_{n2}, \dots, w_{n L_n}), \qquad w_{n\ell} \in \mathcal{V},
\]
where $L_n$ is the length of document $n$ and $\mathcal{V} = \{v_1, \dots, v_V\}$ is the
vocabulary of $V = |\mathcal{V}|$ distinct terms. Each document belongs to a district
$i(n) \in \{1, \dots, I\}$ ($I = 1{,}485$ districts, indexed by NCES LEAID) and a time period
$t(n)$ (calendar year or quarter of the meeting date). Each district has a partisan label
$P_i \in \{\mathsf{R}, \mathsf{D}\}$, equal to $\mathsf{R}$ if its county's 2020 two-party
Republican presidential vote share exceeds one half. To avoid a clash with the corpus
$\mathcal{D}$, we write the parties in sans-serif and the sets of districts as
$\mathcal{I}^{\mathsf{R}}$ and $\mathcal{I}^{\mathsf{D}}$.

\subsection{Bag-of-words representation}\label{sec:m-bow}

We represent each document by its bag-of-words count vector
$\mathbf{c}_n \in \mathbb{N}^V$, so that the corpus is an $N \times V$
document--term matrix. Because partisanship is
measured at the district level, we sum counts within district $i$ and period $t$,
giving $\mathbf{c}_{it}$ with frequencies $\mathbf{q}_{it} = \mathbf{c}_{it}/m_{it}$,
where $m_{it} = \sum_j (c_{it})_j$. Each district thus receives equal weight regardless
of how many meetings it posts, and the partisanship measure and lexicon rates below
are computed from $\mathbf{q}_{it}$.

Following standard practice, we also reweight counts by TF-IDF, which downweights
terms that are common to almost every district and emphasizes distinctive ones. For
district-period $(i,t)$ and term $j$,
\[
  (x_{it})_j = \bigl(1 + \log (c_{it})_j\bigr)\left(\log\frac{1 + M}{1 + \mathrm{df}_j} + 1\right),
\]
where $M$ is the number of district-periods and $\mathrm{df}_j$ is the number that use
term $j$; each row is then normalized to unit length. We use these vectors to find
the terms most characteristic of each period and party, and to compare cosine
similarity between same-party and different-party pairs of districts.

\paragraph{Dictionary of themes.}
To measure specific agendas, we define 17 themes by hand, each a list of single
words and multi-word phrases (for example, ``critical race theory,'' ``book ban''
and ``point of order''). Four themes cover routine school business and serve as
controls; the other 13 cover political and contested themes. Single words are
counted from the document--term matrix and phrases directly from the cleaned text,
and each theme's rate is its mentions per 10,000 words in $\mathbf{c}_{it}$. Each
theme carries a prior about which party discusses it more, which the analysis tests
rather than assumes; partisanship always comes from county vote share. The word
lists are in the code repository.


\subsection{Topic Modeling}\label{sec:m-comm}

If the vocabulary used in school board meetings changes between the periods before and after COVID, it seems plausible that the topics would as well. As such, topic modeling was applied to the dataset with LDA \citep{blei_latent_2003}. For $K$ topics, the topic model proposes
\[ d_i \sim \text{Multinomial}(L_i, \Omega_i), \quad \Omega_i = \sum_{k=1}^K \beta_i(k) A_k \]
\noindent where $A_k \in \mathbb R^p$ is a probability mass function over the $p$ words for topic $k$ so $A_k(j) \geq 0$ and $\sum_j A_k(j)=1$ and $\beta_i \in \mathbb R^K$ is the topic weight vector of document $i$ with $\beta_i(k) \geq 0$ and $\sum_k \beta_i(k)=1$. Thus, each document's word distribution $\Omega_i$ is a mixture of the $K$ topics. To make the model identifiable, LDA assumes that the weights $\beta_i$ are drawn from a Dirichlet distribution so $\beta_i \overset{iid}{\sim} \text{Dirichlet}(\alpha)$. Equivalently, each word $m$ of document $i$ is generated by drawing $\beta_i$, a topic $z_{im} \sim \text{Multinomial}(1, \beta_i)$, and a word $w_{im}|(z_{im}=k) \sim \text{Multinomial}(1, A_k)$. 

We use a symmetric prior $\alpha(k) = 0.1$ and a Dirichlet$(0.01 \cdot \mathbf 1_p)$ prior for each topic. Because the marginal density $\int p(\beta_i; \alpha) p(d_i | \beta_i, A) d\beta_i$ is intractable, we fit the model using variational inference, in which the posterior $(\beta_i, z_{im}| d_i, \alpha, A_k)$ is approximated by factorized Dirichlet and multinomial distributions. We apply the variant of \citet{hoffman2010lda}. Each chunk of our dataset is treated as a document, and a meeting's topic weights are the average of its chunks' weights, weighted by each chunk's number of in-vocabulary words. We hold out 10\% of meetings in a test set for evaluation while the train set is used to estimate vocabulary and topic-word distributions.

\subsubsection{Higher Criticism}\label{sec:m-lda} We test whether mean within-district word shares changed between
2018--2019 and 2022--2023. Starting from the 83,369-recording
HC source sample, we require at least four eligible recordings
spanning at least three distinct calendar quarters in each of
2018, 2019, 2022, and 2023. This yields 21,878 recordings from
273 districts; coverage in 2020 and 2021 is not required.
We retain all eligible recordings in the four comparison years.
For each recording, we calculate word shares within a fixed
12,000-term vocabulary, average recordings within district-year,
and average the two years within each period. Districts receive
equal weight. Support and variance filters leave 11,804 tested
terms.

Let $\mathcal I_{\mathrm{HC}}$ denote the HC district sample,
with $I_{\mathrm{HC}}=|\mathcal I_{\mathrm{HC}}|=273$.
For this analysis, $\mathcal V_{\mathrm{HC}}$ denotes the fixed
12,000-term vocabulary, and $\mathcal J_{\mathrm{HC}}$ indexes
its $m=11{,}804$ tested terms. For district $i$ and vocabulary
term $v_j$, let $\Delta_{ij}$ be the late-minus-early word share,
using the recording and year averages described above. We calculate
\[
\overline{\Delta}_j
=
\frac{1}{I_{\mathrm{HC}}}
\sum_{i\in\mathcal I_{\mathrm{HC}}}\Delta_{ij},
\qquad
Z_j
=
\frac{\sqrt{I_{\mathrm{HC}}}\,
      \overline{\Delta}_j}{s_j},
\]
where $s_j$ is the sample standard deviation of $\Delta_{ij}$
across districts. The two-sided normal-approximation p-value is
$p_j=2\Phi(-|Z_j|)$, where $\Phi$ is the standard normal cumulative
distribution function. The global null is zero mean change for
every tested term. Ordering the p-values as
$p_{(1)}\leq\cdots\leq p_{(m)}$, we calculate
\begin{equation}
\mathrm{HC}^{+}
=
\max\left\{
0,\;
\max_{\substack{1\leq k\leq\lfloor m/2\rfloor\\
                p_{(k)}>1/m}}
\sqrt{m}\,
\frac{k/m-p_{(k)}}{\sqrt{p_{(k)}(1-p_{(k)})}}
\right\}.
\label{eq:hc}
\end{equation}
An empty maximum is set to zero. We also report
$\max_{j\in\mathcal J_{\mathrm{HC}}}|Z_j|$ and
$\sum_{j\in\mathcal J_{\mathrm{HC}}}Z_j^2$.

We calibrate all three global statistics using a Gaussian
district-multiplier bootstrap; implementation, assumptions,
and sensitivity checks are given in
Appendix~\ref{app:hc-calibration}.

\paragraph{Phrase-based follow-up.}
Six exploratory dictionaries measure equity/inclusion, race/racism,
book restrictions, parents' rights, school choice/vouchers, and
closures/reopening. Rates are matches per 10,000 cleaned alphabetic
tokens, including stopwords, with the same recording/year/district
weighting. We report paired-district mean changes and pointwise
95\% t intervals, and apply Holm adjustment across the six
categories separately for each specification. A second
specification removes entire approximately 1,000-word chunks
containing COVID terms or candidate misrecognitions, recalculates
rates on the remaining text, and reapplies the coverage requirements.
Dictionaries were refined after context inspection; these analyses
are exploratory measures of mentions, not stance or causal effects.

\subsubsection{Spectral Clustering}
Having obtained a topic weight vector $\beta_i$ for each of the documents, it becomes of interest to investigate how districts cluster based on the topics they discuss and whether political partisanship varies across clusters. To do this, spectral clustering was performed on a matrix $X \in \mathbb R^{I \times K}$ where $X_{ij}$ is the weight that district $i$ places on topic $j$. Row $i$ of $X$ is equal to $\gamma_i^T$, where the weight vector $\gamma_i \in \mathbb R^K$ is obtained for district $i$ by averaging the weight vectors for the set of documents $\mathcal S_i$ that belong to district $i$: $\gamma_i = \frac{1}{|\mathcal S_i|} \sum_{j \in \mathcal S_i} \beta_j$.

$X$ is standardized by removing the mean of each column and dividing by the standard deviation within each year to obtain $X^*$, and singular value decomposition is performed so $X^* = UDV^T$, where orthogonal matrices $U$ and $V$ contain the left and right singular vectors (respectively) and diagonal $D$ the singular values (in non-increasing order). Let $X$ have rank $r$. Because standardizing removes one degree of freedom, $X^*$ has rank $r-1$, so with noise, the rows of $U$ roughly divide into $r$ clusters. After determining the rank of $X^*$ from a plot of its singular values, $k$-means clustering into $r$ clusters is performed on the first $r-1$ left singular vectors of $X^*$ and examined for prominent topics and political partisanship per cluster.


\subsection{LLM-Based Methods}\label{sec:m-llm}

One important limitation of word-count-based methods is that they do not reveal which side a speaker takes. For example, ``mask mandate'' is counted the same whether a parent demands or objects to one, and ``equity'' is the same in a staff report as in a shouted objection. Topic models share the limitation, since a topic is a distribution over words rather than a position. We therefore use a large language model, as in \citet{asirvatham_gpt_2026}, to answer focused positional questions about every passage, so that qualitative judgments of the kind a human coder would make are produced at scale and can be compared across districts and years.

GABRIEL \citep{gabriel} is a thin wrapper around a model API: its \texttt{rate} function sends each text with a dictionary of attribute definitions in a single prompt and parses the reply into a 0--100 score per attribute. Nothing is learned from our data; the model applies its background knowledge of U.S.\ political discourse to the rubric we supply, so the rubric wording is the entire specification of the measure. \citet{asirvatham_gpt_2026} report that such scores match human coders across many labeled tasks and are robust to rewording; \citet{ludwig_llm_2025} note that a model-generated score used as a regression variable should be validated, since a different model or prompt can change estimates. We treat the scores as descriptive measurements and validate them by repeat runs, agreement between two models of different size, and reading (\cref{app:llm}).

\paragraph{Sample and attributes.}
Because the model reads passages in context, we work from the raw captions rather than the cleaned corpus. We keep 348 districts with at least two eligible recordings in every year of 2018--19 and 2021--23, sample three recordings per district-year (six in 2020, 2021 and 2024, added to locate the timing of any change and test its persistence), and from each of the 10,167 meetings take one random window of about 1,500 words plus up to two windows centered on contested-topic lexicon terms. Random windows measure typical meeting content; flagged windows measure how contested topics are argued once raised, and the two are never pooled. District, place and personal names are replaced with placeholders so the model cannot infer a district's politics from its name. Each excerpt is scored on six attributes in one call: hostility; share of discussion on culture-war issues; parental-rights framing; contention over COVID policy; and the strength with which speakers argue the progressive side, and the conservative side, of a contested dispute. The position attributes count only positions argued in a disagreement, so routine descriptions of district programs score zero. To prevent tuning the attributes to the result, districts are split in half by a deterministic hash of their NCES identifier; wording was developed on a 200-excerpt pilot from split A, revised once, and frozen, and split B is held out to confirm the results. All 28,374 excerpts are scored by \texttt{gpt-6-luna}; a separate pass scores the share of each excerpt that is public comment rather than board or staff discussion. Sample construction, attribute definitions, name removal and validation details are in \cref{app:llm}.

\paragraph{Estimation.}
We average scores to the district-year and, for each attribute, estimate
\begin{equation}
  y_{it}=\alpha_i+\gamma_t+\sum_{s\neq 2019}\beta_s\,
  \mathbf{1}\{t=s\}\,G_i+\mathbf{x}_i'\boldsymbol{\delta}_t+\varepsilon_{it},
  \label{eq:es}
\end{equation}
where $y_{it}$ is the mean score in district $i$ and year $t$, $G_i$ is the county 2020 two-party Republican share in tens of percentage points, $\alpha_i$ and $\gamma_t$ are district and year effects, $\mathbf{x}_i$ contains locale indicators interacted with year, and standard errors are clustered by district. Each $\beta_s$ measures how the association between county partisanship and the score in year $s$ differs from 2019; district effects restrict identification to within-district change. We re-estimate \cref{eq:es} at half-year resolution for 2020--21, separately for excerpts above and below 50 on the public-comment share, and with a single post-period indicator (2021--24) for each sample half. We use \texttt{statsmodels} 0.15 with package defaults; the design is descriptive, not causal.
% =====================================================================
\section{Results}\label{sec:results}

\subsection{Bag-of-Words}\label{sec:r-bow}

We count mentions of the 17 dictionary themes (\cref{sec:m-bow}). Four cover
mainstream school business (budget, facilities, athletics and academics) and serve
as controls. The other 13 cover political and contested themes: parental rights, book content, critical race theory and race,
gender and LGBTQ issues, school choice, COVID mandates, equity and social-emotional
learning, guns and safety, unions and pay, immigration, religion, and two measures of
conflict, in-meeting disruption and legal disputes.

\Cref{fig:lexicon-gap} documents partisan differences in how often these themes are
mentioned in 2022--23. The largest gaps are in professional and local vocabulary
rather than culture-war topics: Democratic-county districts use far more equity and
social-emotional-learning language ($t=-17.5$), along with more union, pay and
academic terms, while Republican-county districts talk more about athletics
($t=9.8$) and facilities ($t=6.8$). By contrast, critical race theory, book content,
parental rights and school choice do not differ detectably ($|t|<1.3$).

\Cref{fig:lexicon-trend} shows the time trend for six key themes of education
politics. COVID mandates and critical race theory spiked in both groups in mid-2021
and faded by 2022. Meeting conflict followed the same path: after dipping in 2020,
when many meetings were virtual, it peaked in the third quarter of 2021 and has since
fallen back to about its pre-pandemic level. Book content, parental rights and gender
issues rose later and stayed above their 2018--19 levels through 2023, in Republican-
and Democratic-county districts alike. \Cref{fig:conflict-state} shows the geography:
conflict language rose in many states between 2018--19 and 2022--23, but most in New
Jersey, Oklahoma and Florida, which reach about 1.8--1.9 mentions per 10,000 words.

\subsection{Topic Modeling}\label{sec:r-comm}
The outputs from LDA were obtained for \textcolor{red}{UPDATE THIS: $K=10,20,30$ topics, with $K=20$ topics being selected because it achieved the lowest perplexity and highest coherence scores while remaining interpretable (see Appendix for details).} Some of the most prominent topics across the entire corpus include informal discussion, voting, budget and finance, and scheduling. Table \ref{tab:lda_topics} displays key words and labels for each topic. The trend of these topics over time is examined by averaging the estimated topic weights across meetings within the same district and then across districts within the same year. This analysis is restricted to the 363 districts that were present across all 6 years. Figure \ref{fig:topic_trends} displays the trends for each topic, highlighting in color the topics that experienced the most change. Notably, the COVID/Health and parental concerns topics peaked during COVID years (2020-2021) while student activities dropped. In both cases, the topics roughly returned to their pre-COVID levels of prevalence in 2022 and 2023. The topic regarding voting interestingly increased over time, including post-COVID.

\subsubsection{Higher Criticism}
Recorded word shares differ between 2018--2019 and 2022--2023.
HC$^{+}$ equals 859.85, compared with a 95th-percentile bootstrap
threshold of 11.34. All three global tests return Monte Carlo
p-values of 0.0005, the reporting limit. After excluding specified
COVID-related, number-format, and suspected artifact words,
11,745 words remain and HC$^{+}$ equals 819.30; all three tests
again reach the reporting limit.

The phrase analysis finds increases in equity/inclusion,
race/racism, book-restriction, and parents' rights mentions
(Table~\ref{tab:hc-themes}). All four survive Holm adjustment
and persist outside COVID-containing chunks, retaining 273
districts and 21,815 recordings. School-choice/voucher and
closure/reopening measures show no detectable change. 

\subsubsection{Spectral Clustering}
The scree plot warranted a selection of 3 clusters, as shown in Figure \ref{fig:lda_scree}. Figure \ref{fig:lda_clustering} shows the clustering results, along with the average GOP share per cluster. Although the districts do not separate into discrete clusters, the clusters still reveal notable differences in political partisanship, with one leaning Democrat, one Republican, and one neutral. This separation persists significantly across the years, with a mild polarization during COVID (Figure \ref{fig:lda_cluster_gop_time}). It is notable, however, that the clusters also tend to align with locale, as the Democratic-leaning cluster is composed of 41.5\% urban districts while the Republican-leaning cluster is 25.3\% rural. As an example, Figure \ref{fig:lda_gop_map} plots the clustering results for districts in year 2021 on top of a heatmap for county GOP share in 2020. Because geographic proximity to a city tends to correlate with political partisanship, ensuring that these clusters were not simply recovering locale becomes of interest. To this end, the share of variance in count GOP vote share explained by cluster membership was compared with its distribution when cluster labels were shuffled within locale (city, suburb, town, rural, for 2,000 permutations). In every year from 2018 to 2023, clusters explained 12-17\% of the variance in partisanship, significantly more than expected from locale alone.

\subsection{LLM-Based Methods}\label{sec:r-llm}

\paragraph{Conservative positions were rarely argued before the pandemic, appeared everywhere in 2021, and most in Republican-leaning counties.}
\Cref{fig:llm-headline} shows yearly means in random excerpts. In 2018--19 the conservative-positions score averages 0.3 in both Republican- and Democratic-county districts, and 2020 is indistinguishable from the pre-period. In 2021 the score rises to 4.3 in Democratic and 6.8 in Republican counties, declines thereafter, and the two groups meet again at 0.9 in 2024. Parental-rights framing follows the same path. In the event study (\cref{tab:llm-es}), ten points of county Republican share are associated with 1.14 more points of conservative positions in 2021 (95\% CI 0.64--1.63) and 0.54 in 2023 (0.11--0.96), relative to 2019; the 2018, 2020, 2022 and 2024 coefficients all include zero. Collapsing the post-2020 years into a single indicator, the coefficient is $+0.45$ ($p=.002$) overall and $+0.55$ ($p=.014$) in the held-out split B; for parental-rights framing it is $+0.45$ ($p=.001$) and $+0.51$ ($p=.004$). The gap is not permanent: it opens in 2021, reappears in 2023 as disputes move from masks to library books, and has closed by 2024.

\paragraph{The disagreement and conflict mainly dates to the second half of 2021 and is concentrated in public comment.}
The half-year coefficients in \cref{tab:llm-es} are near zero in both halves of 2020 and small in the first half of 2021. The gap in conservative positions appears in 2021H2 ($+1.88$, CI 1.05--2.71), the opening of the 2021--22 school year, when mask mandates returned and curriculum disputes peaked; contention over COVID policy spikes in the same half-year ($+2.13$). The change did not occur while schools were closed. Splitting by who is speaking, excerpts that are mostly public comment show a 2021 coefficient of $+12.1$ (CI 5.3--18.8) for conservative positions, against $+0.42$ (0.07--0.78) for excerpts that are mostly board or staff discussion, which look alike across Republican- and Democratic-county districts in every year. Public comment is only 6--10\% of random excerpts, so these estimates are imprecise and only 2021 excludes zero. 

To make these results more tangibile, we pull quotes from Indian River County, Florida (Republican share 0.61) to illustrate the sequence in verbatim: "please make masks optional and allow parents to make this decision for their children" (September 2020); "you will have opened the floodgates of an indoctrination" (May 2021, on critical race theory); "it's not too late to rescind the vote and follow the laws and statutes that protect children" (March 2022, on library books). Most of the district's 34 verified quotes are public comment (\cref{app:llm}). Indian River County was also a hotbed for Moms for Liberty, a far-right parental rights movement founded in 2021 \citep{splc_moms_2023}.  

\paragraph{Hostility rose in both groups}
Random-excerpt hostility rose from 2019 to 2021 in Democratic (3.7 to 5.2) and Republican (2.5 to 5.2) counties and remained elevated. Republican-county districts were the less hostile group before the pandemic (2019 gap $-1.15$, CI $-2.12$ to $-0.21$) and converged; the post-period coefficient ($+0.39$, $p=.007$) does not replicate in split B ($+0.25$, $p=.16$). Progressive positions show no partisan change. In flagged excerpts (\cref{fig:llm-topics}), hostility rose in both groups on every topic except equity, and both groups moved in the same direction on most topics; books and race in the curriculum are the exceptions, with Republican-county meetings well to the conservative side after 2020.

\paragraph{Limitations.}
There are several important limitations to these results, including the fact we use county vote share as proxies for district partisanship, public-comment share is inferred from caption text, and scores are validated by cross-model agreement and reading rather than human coding. Within those limits the results agree with \cref{sec:r-bow}: culture-war topics rose everywhere at once, and for roughly two years, through public comment rather than board discussion, the positions argued diverged by county politics and then converged. 

% =====================================================================
\section{Discussion}\label{sec:discussion}

We make three main contributions. First, we provide a systematic, multi-method
text analysis of DistrictView, a newly released corpus of school board meetings,
complementing work that measures conflict in local politics
\citep{holman_et_al_2026}. 

Second, we provide a broad, quantitative analysis of how education politics changed
around the pandemic, complementing qualitative evidence. We show that contested topics such as critical race theory, book content
and parental rights became salient rhetoric in Republican- and Democratic-leaning districts
at the same time, and hostility rose in both. Partisan differences in language
grew only slightly and mostly reverted by 2023. School closures account for part
of that partisan divergence, but not for the broader rise in conflict, which was
shared universally across districts. The clearest partisan split came in late 2021, when
conservative positions argued in meetings rose more in Republican-leaning
counties, mostly during public comment.  Its timing coincides with the rise of
parents'-rights organizations such as Moms for Liberty, which future work should consider. 

Third, we demonstrate the comparative value of text analysis methods for quantiative social science tasks.  Simple bag-of-words
measures are surprisingly informative about which topics boards discuss and which
words distinguish partisan contexts, but they record mentions rather than
positions. Higher Criticism provides a formal hypothesis test allowing us to confirm that meeting vocabulary changed between 2018--19 and 2022--23 across many weak word-level signals. Topic
models summarize the meeting agenda but are dominated by routine governance.
LLM scoring captures the positions speakers argue, which word-based methods miss,
at the cost of depending on the wording of the scoring rubric and the model used.

Our findings have several limitations. The sample includes only districts that
post meetings on YouTube and is not representative of all U.S.\ districts. County
vote share is a proxy for district partisanship. Automatic captions contain
transcription errors and almost never identify speakers. Our word-based measures
rely mostly on single words, and many topics capture procedural rather than
political content. Future work could combine LLM-based measurement with formal
hypothesis tests, use district-level election returns, or link changes in meeting
content to the formation of local parents'-rights groups.

% =====================================================================

{\small
\bibliography{references}
}

% =====================================================================

\clearpage
\appendix
% =====================================================================
%  Appendix (input from main.tex after \appendix; no preamble here)
% =====================================================================
\section*{Figures}

\begin{figure}[htbp]
  \centering
  \includegraphics[width=0.7\textwidth]{figures/workflow.png}
  \caption{Project workflow. DistrictView transcripts are pre-processed and
  represented as bag of words (TF-IDF) and LDA topic models. Conflict clusters
  hand-picked from the bag-of-words output feed the LLM scoring, and the topic
  model feeds the higher-criticism test and the spectral clustering of districts.}
  \label{fig:workflow}
\end{figure}

\begin{figure}[ht]
  \centering
  \includegraphics[width=0.75\textwidth]{figures/lexicon_category_gap.png}
  \caption{Partisan differences in contested-topic and control vocabulary, 2022--23.
  Bars show Welch $t$-statistics for the difference in district mention rates per
  10,000 words between Republican- and Democratic-county districts (positive = more
  common in Republican counties). Categories with $|t| > 3$ are colored; dotted lines
  mark $\pm 3$. Categories prefixed \texttt{ctrl} are non-contested controls. No
  locale controls.}
  \label{fig:lexicon-gap}
\end{figure}

\begin{figure}[ht]
  \centering
  \includegraphics[width=\textwidth]{figures/gabriel_headline.png}
  \caption{LLM-scored conservative positions and parental-rights framing in random excerpts, 2018--2024. Top: yearly means for districts in Republican- and Democratic-leaning counties, 95\% bootstrap bands over districts. Bottom: Republican minus Democratic. Grey marks 2020, when many meetings were virtual. 348 districts, 16,367 random excerpts.}
  \label{fig:llm-headline}
\end{figure}

\begin{figure}[t]
\centering
\begin{subfigure}[b]{0.49\linewidth}
    \centering
    \includegraphics[width=0.9\linewidth]{figures/lda_topic_trends.png}
    \caption{Topic prevalence over time}
    \label{fig:topic_trends}
\end{subfigure}
\hfill
\begin{subfigure}[b]{0.49\linewidth}
    \centering
    \includegraphics[width=0.9\linewidth]{figures/cluster_gop_2021.png}
    \caption{District clusters and county GOP share, 2021}
    \label{fig:cluster_gop_2021}
\end{subfigure}
\caption{(a) Mean topic weight by year, averaged across districts, with 95\% district-bootstrap intervals. (b) District topic clusters in 2021, plotted over 2020 county GOP vote share.}
\label{fig:trends_and_map}
\end{figure}
\begin{figure}[ht]
  \centering
  \includegraphics[width=\textwidth]{figures/lexicon_category_trend.png}
  \caption{Mentions of contested topics over time. Quarterly mean of district mention
  rates per 10,000 words, Republican-county (red) and Democratic-county (blue)
  districts, 2018--2023. Vertical scales differ across panels. Mentions count words
  used, not the position taken.}
  \label{fig:lexicon-trend}
\end{figure}

\begin{figure}[ht]
  \centering
  \includegraphics[width=\textwidth]{figures/conflict_state_by_period.png}
  \caption{Conflict language by state and period. Each tile shows the state mean of
  district mention rates per 10,000 words for meeting-conflict terms (e.g.\ ``out of
  order'', ``point of order'') and legal-conflict terms (e.g.\ lawsuit, resignation,
  censure). States with fewer than five sample districts are grey. Captions carry no
  tone, so this measures conflict that is said, not heard.}
  \label{fig:conflict-state}
\end{figure}
\begin{figure}[H]
  \centering
  \includegraphics[width=\textwidth]{figures/gabriel_topic_period.png}
  \caption{Flagged excerpts: mean stance (conservative minus progressive) and hostility by contested topic, 2018--19 (open) and 2021--24 (filled), Republican- and Democratic-county districts, 95\% bootstrap intervals over districts.}
  \label{fig:llm-topics}
\end{figure}

\begin{figure}[H]
  \centering
  \includegraphics[width=\textwidth]{figures/gabriel_event_study_by_speaker.png}
  \caption{Event study by who is speaking: coefficients on county Republican share $\times$ year (2019 reference) for excerpts that are mostly public comment (share $\ge 50$) and mostly board or staff discussion, random excerpts, district and year fixed effects, 95\% CIs clustered by district. Public-comment cells hold 55--199 excerpts per year.}
  \label{fig:llm-speaker}
\end{figure}

\begin{figure}[H]
  \centering
  \includegraphics[width=\textwidth]{figures/gabriel_event_study.png}
  \caption{Event study for all seven LLM scores (the six attributes and stance), random excerpts, 2019 reference, district and year fixed effects, locale$\times$year controls, 95\% CIs clustered by district.}
  \label{fig:llm-es-all}
\end{figure}

\begin{figure}[H]
  \centering
  \includegraphics[width=\textwidth]{figures/gabriel_change_map.png}
  \caption{Change in random-excerpt hostility (top) and conservative positions (bottom) from 2018--19 to 2021--24, state means over sampled districts; states with fewer than five districts are grey, and each tile shows the number of districts.}
  \label{fig:llm-map}
\end{figure}

\section*{Tables}

\begin{table}[H]
\centering
\caption{Event-study estimates, \cref{eq:es}: coefficients per 10 percentage points of county 2020 Republican vote share, relative to 2019, with district and year fixed effects and locale$\times$year controls; 95\% confidence intervals clustered by district in brackets. Random excerpts. Lower panel: half-year resolution for 2020--21.}
\label{tab:llm-es}
\footnotesize
\setlength{\tabcolsep}{4pt}
\resizebox{\textwidth}{!}{\input{tables/gabriel_event_study}}
\end{table}

\begin{table}[t]
\centering
\small
\caption{Topics from LDA, 2018--2023. Top words are the eight highest-probability words in each topic $\hat A_k$. Share is the mean topic weight
$\hat w_i(k)$, averaged over meetings within each district and then across districts.}
\label{tab:topics}
\begin{tabularx}{\linewidth}{@{}l l X r@{}}
\toprule
 & Label & Top words & Share (\%) \\
\midrule
\multicolumn{4}{@{}l}{\textit{Panel A: Substantive topics}} \\
T14 & Budget \& finance        & year, budget, dollars, million, money, fund, state, thousand & 8.5 \\
T15 & Calendar \& scheduling   & day, time, days, week, th, think, teachers, staff & 7.3 \\
T5  & Facilities \& construction & building, project, new, projects, construction, work, high, cost & 5.7 \\
T4  & Strategic planning       & work, plan, really, district, community, data, process, year & 5.3 \\
T12 & Health, safety \& COVID  & children, parents, health, students, people, kids, schools, no & 5.0 \\
T2  & Achievement data         & students, year, grade, high, data, look, level, student & 4.9 \\
T9  & Student services         & students, program, schools, staff, services, support, education, district & 4.9 \\
T10 & Curriculum \& instruction & students, teachers, learning, teacher, curriculum, really, classroom, student & 4.8 \\
T13 & Athletics \& activities  & high, team, bus, year, sports, state, students, field & 3.7 \\
T1  & Municipal \& land use    & city, street, property, town, no, area, parking, building & 1.6 \\
\addlinespace
\multicolumn{4}{@{}l}{\textit{Panel B: Procedural and conversational topics}} \\
T6  & Informal discussion      & really, kids, things, lot, did, think, great, kind & 12.4 \\
T3  & Deliberation             & think, say, need, mean, make, people, board, question & 10.9 \\
T11 & Ceremony \& recognition  & board, students, community, district, mr, dr, evening, high & 8.7 \\
T8  & Meeting procedure        & board, meeting, policy, public, committee, members, agenda, th & 8.3 \\
T7  & Motions \& voting        & yes, motion, mr, second, approve, aye, favor, item & 8.1 \\
\bottomrule
\end{tabularx}
\end{table}
\begin{table}[ht]
\centering
\caption{Summary of the analysis corpus.}
\label{tab:corpus}
\small
\begin{tabular}{lrrrr}
\toprule
Group & Districts & Meetings & Words (millions) & Period \\
\midrule
\multicolumn{5}{l}{\textit{Panel A: by county partisanship}} \\
Republican county & 703 & 31,260 & 380.4 & 2018--23 \\
Democratic county & 779 & 48,517 & 743.9 & 2018--23 \\
Unclassified (Alaska) & 3 & 286 & 4.9 & 2018--23 \\
\midrule
\multicolumn{5}{l}{\textit{Panel B: by period}} \\
Pre-COVID & 545 & 12,866 & 161.1 & 2018--19 \\
COVID & 1,277 & 30,516 & 453.3 & 2020--21 \\
Post-COVID & 1,339 & 36,681 & 514.8 & 2022--23 \\
\midrule
Total & 1,485 & 80,063 & 1,129.1 & 2018--23 \\
\bottomrule
\end{tabular}

\vspace{0.5em}
\parbox{0.9\linewidth}{\footnotesize \textit{Notes:} A meeting is one YouTube video
that passes the sample filters. Words are counted in the cleaned transcripts before
stopword and vocabulary cuts. Partisanship is the district-office county's 2020
two-party Republican vote share, split at 0.5. Districts in Panel B do not sum to
the total because most appear in more than one period.}
\end{table}

\begin{table}[htbp]
\centering
\caption{Changes in policy-related mentions per 10,000 cleaned
alphabetic tokens, 2018--2019 to 2022--2023.
Brackets give pointwise 95\% confidence intervals.}
\label{tab:hc-themes}
\small
\begin{tabular}{lrr}
\toprule
Category & All text & Outside COVID chunks \\
\midrule
Equity/inclusion
 & \shortstack{0.778\\{}[0.513, 1.044]}
 & \shortstack{0.816\\{}[0.545, 1.086]} \\
Race/racism
 & \shortstack{0.220\\{}[0.152, 0.287]}
 & \shortstack{0.234\\{}[0.162, 0.306]} \\
Book restrictions
 & \shortstack{0.0557\\{}[0.0381, 0.0733]}
 & \shortstack{0.0571\\{}[0.0391, 0.0752]} \\
Parents' rights
 & \shortstack{0.0965\\{}[0.0649, 0.1280]}
 & \shortstack{0.0987\\{}[0.0665, 0.1308]} \\
School choice/vouchers
 & \shortstack{$-0.006$\\{}[$-0.051$, 0.039]}
 & \shortstack{$-0.012$\\{}[$-0.053$, 0.030]} \\
Closures/reopening
 & \shortstack{0.017\\{}[$-0.002$, 0.035]}
 & \shortstack{0.006\\{}[$-0.013$, 0.025]} \\
\bottomrule
\end{tabular}
\end{table}

\section{Corpus construction and cleaning}\label{app:cleaning}

All analyses start from one cleaned corpus (\cref{sec:cleaning}). This appendix
gives the sample counts at each step, the normalization and screening rules, and
how the cleaned text is converted to counts for the bag-of-words and Higher
Criticism (HC) analyses.

\subsection{Sample construction}
We select records with a nonmissing meeting date in 2018--2023, giving 110,777
recordings. We exclude 9,914 transcript placeholders, 1,216 recordings whose
transcripts are missing from the archives, 1,627 transcripts with fewer than 200
word tokens after normalization, and 94 exact duplicates, leaving 97,926
recordings. Duplicate transcripts are identified by identical normalized word
sequences; the first encountered copy is retained. The school-board eligibility
screen below leaves 83,369 recordings, and requiring at least 1,000 words leaves
the final 80,063 meetings from 1,485 districts. \todo{Confirm that the 83,369 $\to$ 80,063 step is the 1,000-word floor alone, and that HC uses 83,369 (200-word floor) or the 80,063 sample; state which.}

\subsection{Text normalization}
We lowercase the text, convert curly apostrophes to straight apostrophes, and
remove bracketed annotations other than timestamps. We extract sequences of
letters a--z, allowing an internal apostrophe; numerals, punctuation and caption
timestamps do not enter the resulting word sequence.

\subsection{School-board eligibility screen}
Titles must either explicitly identify school-board proceedings or contain a
governance term supported by channel or transcript evidence. Channel support
requires at least ten explicitly titled school-board recordings, constituting at
least half the channel's screened recordings, together with at least three
education-related word occurrences covering two distinct terms in the first 600
transcript words. Alternatively, the opening text must explicitly identify a
school governing body. We exclude recordings flagged as possible clips or campaign
events, other public-body meetings, joint-body meetings, or unresolved cases. The
primary eligible sample contains 83,369 recordings: 70,748 with explicit titles and
12,621 with corroborating evidence. Sampled text reviews supplement these rules;
eligibility is not manually verified for every recording.

\subsection{Bag-of-words representation}
Tokens are counted per meeting after removing stop words. \todo{Fill in from the BoW pipeline: stopword list, minimum document frequency (20 meetings), vocabulary cuts (digits, terms under three characters, possessives merged, procedural words, terms used in fewer than 25 districts), and the multi-word phrase counts for the lexicon.} Counts are then summed to district-quarter and district-period for the partisanship and lexicon measures (\cref{sec:r-bow}), and reweighted by TF-IDF for the distinctive-term and similarity analyses.

\subsection{Higher Criticism representation}
HC uses a different count representation of the same cleaned text. A fixed
dictionary expands common contractions and informal forms, such as ``can't'' to
``cannot'' and ``gonna'' to ``going to,'' and fixed suffix rules expand forms such
as ``n't'' to ``not'' and ``'ve'' to ``have.'' We divide each transcript into
consecutive 1,000-word chunks; a final chunk shorter than 200 words is merged into
the preceding chunk. We reuse the vocabulary fitted for the original LDA analysis,
using training chunks only. We remove scikit-learn's English stopwords, except
``no,'' ``not,'' ``nor,'' and ``never,'' and add a fixed list of fillers and
contraction fragments. Candidate terms contain at least two letters and must appear
in at least 25 training chunks and no more than 80\% of training chunks. We retain
the 12,000 eligible terms with the highest total training counts, and counts
exclude all other terms. The matched-district restrictions and additional word
filters are in \cref{sec:m-lda}.

\section{LLM measurement}\label{app:llm}


\paragraph{Sample.}
Eligible recordings pass the screen of \cref{app:cleaning}, have at least 3,000 words, and are unique within district and date (the longest is kept). 349 districts have at least two eligible recordings in every year of 2018, 2019, 2021, 2022 and 2023; 348 have vote share and locale. The 12 strata (Republican-share tercile $\times$ city, suburb, town, rural) are uneven because Democratic rural (5) and Republican urban (11) districts are rare; the sample over-represents suburbs and Democratic-leaning counties. Recordings sampled per year, 2018--2024: 1,032, 1,041, 2,036, 2,060, 1,042, 1,040 and 1,916; half-years in 2020 and 2021 receive three each where available, and 16 districts have no eligible 2024 recording.

\paragraph{Excerpts.}
Random windows have median length 1,495 words. Flagged windows are centered on clusters of at least two lexicon hits, with rarer categories (CRT, parental rights, books, gender) preferred over COVID and equity when several apply; counts by category (post/pre, original sample): COVID 1,496/334; equity/SEL 974/743; gender/LGBTQ 505/271; books 406/147; parental rights 365/81; CRT/race 338/23. A variance decomposition of random-excerpt conservative positions gave 72.7\% within district-year, 21.6\% between years within district and 5.7\% between districts, so two further random windows per meeting were added for 2021 and 2023, narrowing the confidence intervals of the Republican--Democratic gap in those years by 17--31\%. Scored excerpts: 16,367 random and 12,007 flagged; 28,374 in total.

\paragraph{Name removal and leakage.}
District, city, county, state and school names from the metadata, the 50 state names, titles followed by up to two name tokens, phrases ending in ``High School,'' ``Elementary,'' ``School District'' and similar, and spaCy \texttt{en\_core\_web\_sm} \texttt{PERSON}/\texttt{GPE} entities are replaced by \texttt{[NAME]} or \texttt{[PLACE]}; ``DEI'' and ``gender'' are protected from the tagger, and ``I'm X'' is deliberately not stripped. Removal deletes 0.1\% of lexicon hits. About 8\% of excerpts retain a detectable place or school name; the stance--partisanship gradient is present in excerpts without one (random $+1.16$ per 10 pp, $p<.001$) and weaker where one remains ($+0.59$, $p=.05$), so leakage does not generate the result.

\paragraph{Attribute definitions (version 3).}
\emph{Hostility}: personal attacks or insults, shouting or interrupting, accusations of lying, corruption or bad faith, threats, calls for officials to resign or speakers to be removed; polite disagreement is low. \emph{Culture-war share}: share of the excerpt on race in the curriculum, library or classroom content, gender identity, parental rights, masks or vaccines; routine business counts as none. \emph{Parental-rights framing}: how strongly issues are framed as parents' authority (right to know, consent, opt out) rather than the district's. \emph{COVID contention}: disagreement over masks, closures or remote learning, vaccines, quarantines, testing; neutral updates are zero. \emph{Progressive / conservative positions}: how strongly speakers argue that side of a contested dispute, with examples (defending books against removal, DEI or LGBTQ-inclusive policy, mask requirements; demanding removal of books, opposing DEI, asserting parental authority over curriculum, opposing CRT, opposing mask requirements); positions argued, not topics discussed; routine descriptions of district programs score zero. Version 2 lacked the last clause: in the pilot it scored a staff equity-plan presentation 57 on progressive positions, and random excerpts averaged 4.4 progressive against 0.3 conservative. Version 3 was re-piloted on the same 200 split-A excerpts and frozen. Every wording tried is logged in \texttt{source/analysis/gabriel/attribute\_log.md}.

\paragraph{Validation.}
Scoring used \texttt{gpt-6-luna} with reasoning disabled, JSON output and 10 parallel calls; total API cost was \$16.93. Pilot run-to-run correlations on 200 excerpts: 0.97, 0.98, 0.97, 0.99, 0.96, 0.99 for the six attributes in the order above, with mean absolute differences of 0.6--3.3 points. Re-scoring 1,000 excerpts (500 random, 500 flagged) with \texttt{gpt-6.1-terra} gives correlations of 0.95, 0.93, 0.95, 0.98, 0.85, 0.97; Terra scores progressive positions 2.4 points lower on average, applying the routine-description rule more strictly, with the largest disagreement in random excerpts ($r=0.70$) where most scores are near zero. We therefore lead with conservative positions ($r=0.97$) and report stance (conservative minus progressive) only as a secondary measure. Post-period regressions re-estimated on the 1,000 shared excerpts give similar coefficients with either model's scores.

\paragraph{Public-comment share.}
A one-attribute pass over the 22,174 original and extension excerpts asked what share of each is public comment by members of the public, as opposed to board or staff discussion, using cues such as ``state your name and address,'' ``three minutes'' and ``next speaker'' (median 0, mean 12.3). Excerpts with share $\ge 50$ are 5--10\% of random excerpts by year (9.7\% in 2021; 199 excerpts from 149 districts). The extra 2021 and 2023 random windows were not scored for public comment and are excluded from the by-speaker split (\cref{fig:llm-speaker}).

\paragraph{Quotes.}
For Indian River (FL, Republican share 0.61) and Radnor Township (PA, 0.36), chosen from the three largest pre-to-post risers in each group, the 25 highest-scoring excerpts on the relevant position attribute and the 10 highest on hostility were passed in raw form to \texttt{gpt-6.1-terra}, which returned the single sentence most strongly arguing the position. A quote is kept only if it is an exact substring of the raw window (34 of 35 and 26 of 35 kept; failures are excerpts with no sentence returned). Private individuals' names are replaced afterwards; caption errors are preserved. Full lists are in \texttt{output/tables/quotes\_1200930.csv} and \texttt{quotes\_4219920.csv} in the project repository.



\end{document}

